\documentclass[stu]{apa7}

\usepackage[spanish]{babel}
\usepackage[utf8]{inputenc}
\usepackage{hyperref}
\usepackage{enumitem}
\usepackage{geometry}
\usepackage{times}
\usepackage{array}
\usepackage{longtable}
\usepackage{booktabs}
\usepackage{colortbl}
\usepackage{xcolor}
\usepackage{tabularx}

\geometry{margin=1in}

\title{Análisis Documental --- Sistema de información web contable sobre recibos de retención y movimientos de caja de la empresa minera MINCH SRL.}
\author{...}
\affiliation{...}
\course{...}
\professor{...}
\duedate{...}

\setlength{\parindent}{0.5in}

\newcolumntype{L}[1]{>{\raggedright\arraybackslash}p{#1}}
\newcolumntype{C}[1]{>{\centering\arraybackslash}p{#1}}

\begin{document}

\maketitle

\section{Matriz del análisis documental}

\begin{longtable}{|L{4cm}|L{6cm}|L{5cm}|}
\hline
\rowcolor[gray]{0.9}
\textbf{Documento} & \textbf{Descripción} & \textbf{Área donde apoyará} \\
\hline
\endhead
Formulario de retención de impuestos (RC-IVA, IUE, IT) & Documento interno donde se registra el cálculo de retenciones a proveedores sin NIT, aplicando las tasas impositivas según el tipo de bien o servicio (Servicios/Bienes). & Módulo de retenciones \\
\hline
Libro de Excel de movimientos de caja & Archivo digital donde se anotan manualmente los ingresos y egresos de cada cuenta bancaria, clasificados por tipo de movimiento. & Módulo de caja \\
\hline
Recibos de retención emitidos & Comprobantes físicos o digitales que se entregan al proveedor como constancia de la retención realizada. & Módulo de retenciones \\
\hline
Extractos bancarios & Reportes emitidos por las entidades financieras (BNB, BCP, BISA, etc.) que muestran el detalle de movimientos de cada cuenta. & Módulo de libro de bancos \\
\hline
Registro de proveedores & Listado de proveedores con datos de identificación (CI, NIT, nombre, teléfono) utilizado para asociar las retenciones y movimientos. & Módulo de retenciones y caja \\
\hline
Registro de clientes & Listado de clientes con datos de identificación, utilizado para estados de cuenta. & Módulo de estados de cuenta \\
\hline
Normativa tributaria boliviana (Ley 843, RC-IVA, IUE, IT) & Marco legal que regula las tasas de retención y las obligaciones tributarias de la empresa ante Impuestos Nacionales. & Módulo de retenciones \\
\hline
Comprobantes de pago (transferencias y cheques) & Documentos que respaldan cada movimiento bancario, identificados por tipo de pago (T = Transferencia, CH = Cheque). & Módulo de libro de bancos \\
\hline
Conciliaciones bancarias mensuales & Hojas de cálculo donde se compara el saldo contable con el saldo bancario para cada cuenta. & Módulo de caja y reportes \\
\hline
Reportes de retención mensuales & Resúmenes elaborados en Excel que consolidan las retenciones realizadas en un mes, clasificadas por tipo de servicio o bien. & Módulo de reportes \\
\hline
\end{longtable}

\newpage

\section{Recolección de información}

\subsection{Cuestionarios}

\subsubsection{Cuestionario Nº 1 --- Contador General}

\textbf{Cargo:} Contador General \\
\textbf{Objetivo:} Identificar las necesidades del módulo de retenciones y cálculo automático de descuentos. \\
\textbf{Instrucciones:} Responda las siguientes preguntas según su experiencia en el área contable.

\begin{enumerate}
    \item ¿Está de acuerdo con la implementación de un sistema web para la gestión de retenciones y movimientos de caja?
    \item ¿Cómo realiza actualmente el cálculo de descuentos en las retenciones (RC-IVA, IUE, IT)?
    \item ¿Cuánto tiempo semanal destina al registro manual de retenciones en Excel?
    \item ¿Qué información considera indispensable que deba mostrar el módulo de retenciones?
    \item ¿Cómo se genera actualmente el número correlativo de cada recibo de retención?
    \item ¿Qué tipo de reportes de retenciones necesita generar mensualmente?
    \item ¿Qué problemas ha detectado en el registro manual de movimientos de caja?
    \item ¿Qué resultados espera del sistema para considerar que su trabajo es más eficiente?
    \item ¿Cuáles son las secciones o unidades comprometidas con el flujo de retenciones?
    \item ¿Qué información debe contener un estado de cuenta para proveedores y clientes?
\end{enumerate}

\subsubsection{Cuestionario Nº 2 --- Auxiliar Contable}

\textbf{Cargo:} Auxiliar Contable \\
\textbf{Objetivo:} Conocer el proceso operativo diario de registro de movimientos y emisión de recibos. \\
\textbf{Instrucciones:} Responda las siguientes preguntas según sus actividades diarias.

\begin{enumerate}
    \item ¿Cuánto tiempo le toma registrar un movimiento de caja manualmente?
    \item ¿Qué datos registra actualmente en cada movimiento de caja (ingreso/egreso)?
    \item ¿Cómo clasifica los movimientos por tipo (débito/crédito) y por cuenta bancaria?
    \item ¿Utiliza algún formato o plantilla para la emisión de recibos de retención?
    \item ¿Qué dificultades enfrenta al buscar un registro de retención anterior?
    \item ¿Cómo se maneja actualmente la numeración de los comprobantes?
    \item ¿Qué herramientas externas (Excel, papel) utiliza para suplir carencias del sistema actual?
    \item ¿Cuántos movimientos de caja registra aproximadamente por día?
    \item ¿Qué tan frecuente encuentra errores o inconsistencias en los registros?
    \item ¿Qué funcionalidades le gustaría que tuviera el sistema para facilitar su trabajo diario?
\end{enumerate}

\subsubsection{Cuestionario Nº 3 --- Gerente/Administrador}

\textbf{Cargo:} Gerente / Administrador \\
\textbf{Objetivo:} Identificar necesidades de reportes y toma de decisiones. \\
\textbf{Instrucciones:} Responda las siguientes preguntas según su rol gerencial.

\begin{enumerate}
    \item ¿Con qué frecuencia necesita consultar el estado financiero de las cuentas bancarias?
    \item ¿Qué tipo de reportes financieros utiliza para la toma de decisiones?
    \item ¿Cómo accede actualmente a la información de retenciones realizadas en un período?
    \item ¿Qué información debe mostrar el dashboard ejecutivo del sistema?
    \item ¿Necesita exportar reportes en PDF o Excel? ¿Con qué frecuencia?
    \item ¿Qué información considera crítica que el sistema deba mostrar en tiempo real?
    \item ¿Cómo se manejan actualmente los estados de cuenta de proveedores y clientes?
    \item ¿Qué beneficios económicos espera de la implementación del sistema?
    \item ¿Qué indicadores de rendimiento (KPI) considera importantes para el área contable?
    \item ¿Cómo prefiere visualizar la información financiera (tablas, gráficos, ambos)?
\end{enumerate}

\newpage

\subsection{Guía de entrevista}

\subsubsection{Guía de Entrevista Nº 1 --- Contador General}

\begin{tabularx}{\textwidth}{|l|X|}
\hline
\rowcolor[gray]{0.9}
\textbf{Campo} & \textbf{Dato} \\
\hline
Nombre & \_\_\_\_\_\_\_\_\_\_\_\_\_\_\_\_\_\_\_\_\_\_\_\_\_\_\_\_\_\_\_\_\_\_\_\_\_\_\_\_\_\_\_\_\_\_\_\_\_\_\_ \\
\hline
Cargo & Contador General \\
\hline
Institución & MINCH SRL. \\
\hline
\end{tabularx}

\vspace{0.3cm}

\begin{enumerate}
    \item ¿Está de acuerdo con la implementación del sistema web para la gestión de retenciones y movimientos de caja, y qué resultados espera para mejorar la eficiencia del área contable?

    \item ¿Qué aspectos son importantes en la implementación del sistema y cómo mejorará el funcionamiento de la empresa?

    \item ¿Cuáles son las unidades o secciones de la empresa comprometidas con el funcionamiento del sistema de retenciones y caja?

    \item ¿Cómo se genera actualmente el proceso de registro de retenciones y qué información se necesita para emitir un recibo de retención?

    \item ¿Cuáles deben ser los resultados que debe generar el sistema en cuanto a la automatización del cálculo de descuentos?

    \item ¿Cómo se realiza el proceso de conciliación bancaria y qué información se necesita del sistema para agilizarlo?

    \item ¿Qué tipos de reportes financieros son críticos para la toma de decisiones gerenciales?

    \item ¿Cómo se maneja actualmente la numeración de los recibos y qué problemas ha ocasionado el método manual?
\end{enumerate}

\vspace{0.3cm}

\textbf{Observaciones:} \_\_\_\_\_\_\_\_\_\_\_\_\_\_\_\_\_\_\_\_\_\_\_\_\_\_\_\_\_\_\_\_\_\_\_\_\_\_\_\_\_\_\_\_\_\_\_\_\_\_\_\_\_\_\_\_\_\_\_\_\_\_\_\_\_\_\_\_\_\_\_\_\_\_\_\_\_\_\_\_\_\_\_\_\_\_\_\_\_\_\_\_\_\_\_

\newpage

\subsubsection{Guía de Entrevista Nº 2 --- Auxiliar Contable}

\begin{tabularx}{\textwidth}{|l|X|}
\hline
\rowcolor[gray]{0.9}
\textbf{Campo} & \textbf{Dato} \\
\hline
Nombre & \_\_\_\_\_\_\_\_\_\_\_\_\_\_\_\_\_\_\_\_\_\_\_\_\_\_\_\_\_\_\_\_\_\_\_\_\_\_\_\_\_\_\_\_\_\_\_\_\_\_\_ \\
\hline
Cargo & Auxiliar Contable \\
\hline
Institución & MINCH SRL. \\
\hline
\end{tabularx}

\vspace{0.3cm}

\begin{enumerate}
    \item ¿Cuáles son las tareas diarias que realiza en el área de contabilidad y cuánto tiempo le demandan?

    \item ¿Qué dificultades enfrenta al registrar los movimientos de caja en Excel?

    \item ¿Cómo se realiza actualmente la emisión de un recibo de retención paso a paso?

    \item ¿Qué información es la que más consulta o recupera de los registros históricos?

    \item ¿Qué cambios sugeriría para reducir el tiempo que dedica a tareas administrativas manuales?

    \item ¿Cómo maneja las diferencias entre los saldos contables y los saldos bancarios?
\end{enumerate}

\vspace{0.3cm}

\textbf{Observaciones:} \_\_\_\_\_\_\_\_\_\_\_\_\_\_\_\_\_\_\_\_\_\_\_\_\_\_\_\_\_\_\_\_\_\_\_\_\_\_\_\_\_\_\_\_\_\_\_\_\_\_\_\_\_\_\_\_\_\_\_\_\_\_\_\_\_\_\_\_\_\_\_\_\_\_\_\_\_\_\_\_\_\_\_\_\_\_\_\_\_\_\_\_\_\_\_

\newpage

\subsubsection{Guía de Entrevista Nº 3 --- Gerente / Administrador}

\begin{tabularx}{\textwidth}{|l|X|}
\hline
\rowcolor[gray]{0.9}
\textbf{Campo} & \textbf{Dato} \\
\hline
Nombre & \_\_\_\_\_\_\_\_\_\_\_\_\_\_\_\_\_\_\_\_\_\_\_\_\_\_\_\_\_\_\_\_\_\_\_\_\_\_\_\_\_\_\_\_\_\_\_\_\_\_\_ \\
\hline
Cargo & Gerente / Administrador \\
\hline
Institución & MINCH SRL. \\
\hline
\end{tabularx}

\vspace{0.3cm}

\begin{enumerate}
    \item ¿Qué información financiera necesita consultar periódicamente para la toma de decisiones?

    \item ¿Cómo accede actualmente a los reportes de retenciones y movimientos de caja?

    \item ¿Qué tan importante es contar con información actualizada en tiempo real sobre los saldos bancarios?

    \item ¿Qué indicadores considera clave para evaluar la salud financiera de la empresa?

    \item ¿Qué formato de reportes prefiere (PDF, Excel, dashboard interactivo) y por qué?

    \item ¿Qué espera del sistema en términos de reducción de costos y mejora de la productividad?
\end{enumerate}

\vspace{0.3cm}

\textbf{Observaciones:} \_\_\_\_\_\_\_\_\_\_\_\_\_\_\_\_\_\_\_\_\_\_\_\_\_\_\_\_\_\_\_\_\_\_\_\_\_\_\_\_\_\_\_\_\_\_\_\_\_\_\_\_\_\_\_\_\_\_\_\_\_\_\_\_\_\_\_\_\_\_\_\_\_\_\_\_\_\_\_\_\_\_\_\_\_\_\_\_\_\_\_\_\_\_\_

\newpage

\subsection{Guía de observación}

\subsubsection{Guía de Observación Nº 1 --- Proceso de retenciones}

\begin{tabularx}{\textwidth}{|l|X|}
\hline
\rowcolor[gray]{0.9}
\textbf{Campo} & \textbf{Dato} \\
\hline
Nombre del observado & \_\_\_\_\_\_\_\_\_\_\_\_\_\_\_\_\_\_\_\_\_\_\_\_\_\_\_\_\_\_\_\_\_\_\_\_\_\_\_\_\_\_\_\_\_\_\_\_\_\_\_ \\
\hline
Cargo & \_\_\_\_\_\_\_\_\_\_\_\_\_\_\_\_\_\_\_\_\_\_\_\_\_\_\_\_\_\_\_\_\_\_\_\_\_\_\_\_\_\_\_\_\_\_\_\_\_\_\_ \\
\hline
Área & Contabilidad \\
\hline
Proceso observado & Registro y emisión de recibos de retención \\
\hline
\end{tabularx}

\vspace{0.3cm}

\begin{enumerate}
    \item ¿Cómo se realiza el proceso de registro de una retención desde que llega la factura del proveedor?
    \item ¿Quiénes intervienen en el proceso de retención (auxiliar, contador, gerente)?
    \item ¿Cuáles son los aspectos más importantes que se deben registrar en cada retención?
    \item ¿En cuánto tiempo se completa el proceso de emisión de un recibo de retención?
    \item ¿Cuáles de las tareas y actividades son manuales y cuáles están automatizadas?
    \item ¿Qué herramientas o archivos se utilizan durante el proceso?
    \item ¿Se identifican errores frecuentes? ¿De qué tipo?
    \item ¿Cómo se almacenan y organizan los recibos de retención emitidos?
    \item ¿Cuánto tiempo toma buscar una retención emitida anteriormente?
\end{enumerate}

\vspace{0.3cm}

\textbf{Tiempo estimado del proceso:} \_\_\_\_\_\_\_\_ minutos \\
\textbf{Fricción tecnológica observada:} \_\_\_\_\_\_\_\_\_\_\_\_\_\_\_\_\_\_\_\_\_\_\_\_\_\_\_\_\_\_\_\_\_\_\_\_\_\_\_\_\_\_\_\_\_\_\_\_\_\_\_ \\
\textbf{Observaciones:} \_\_\_\_\_\_\_\_\_\_\_\_\_\_\_\_\_\_\_\_\_\_\_\_\_\_\_\_\_\_\_\_\_\_\_\_\_\_\_\_\_\_\_\_\_\_\_\_\_\_\_

\newpage

\subsubsection{Guía de Observación Nº 2 --- Movimientos de caja}

\begin{tabularx}{\textwidth}{|l|X|}
\hline
\rowcolor[gray]{0.9}
\textbf{Campo} & \textbf{Dato} \\
\hline
Nombre del observado & \_\_\_\_\_\_\_\_\_\_\_\_\_\_\_\_\_\_\_\_\_\_\_\_\_\_\_\_\_\_\_\_\_\_\_\_\_\_\_\_\_\_\_\_\_\_\_\_\_\_\_ \\
\hline
Cargo & \_\_\_\_\_\_\_\_\_\_\_\_\_\_\_\_\_\_\_\_\_\_\_\_\_\_\_\_\_\_\_\_\_\_\_\_\_\_\_\_\_\_\_\_\_\_\_\_\_\_\_ \\
\hline
Área & Contabilidad \\
\hline
Proceso observado & Registro de movimientos de caja \\
\hline
\end{tabularx}

\vspace{0.3cm}

\begin{enumerate}
    \item ¿Cómo se realiza el proceso de registro de un movimiento de caja (ingreso o egreso)?
    \item ¿Quiénes intervienen en el proceso de registro de movimientos de caja?
    \item ¿Cuáles son los datos mínimos que se registran en cada movimiento?
    \item ¿En cuánto tiempo se completa el registro de un movimiento de caja?
    \item ¿Cuáles de las tareas y actividades son manuales y cuáles automatizadas?
    \item ¿Qué herramientas externas (Excel, calculadora, papel) se utilizan durante el proceso?
    \item ¿Cómo se verifica que los saldos registrados coincidan con los saldos bancarios?
    \item ¿Con qué frecuencia se realizan conciliaciones bancarias?
    \item ¿Cuánto tiempo toma generar un reporte mensual de movimientos de caja?
\end{enumerate}

\vspace{0.3cm}

\textbf{Tiempo estimado del proceso:} \_\_\_\_\_\_\_\_ minutos \\
\textbf{Fricción tecnológica observada:} \_\_\_\_\_\_\_\_\_\_\_\_\_\_\_\_\_\_\_\_\_\_\_\_\_\_\_\_\_\_\_\_\_\_\_\_\_\_\_\_\_\_\_\_\_\_\_\_\_\_\_ \\
\textbf{Observaciones:} \_\_\_\_\_\_\_\_\_\_\_\_\_\_\_\_\_\_\_\_\_\_\_\_\_\_\_\_\_\_\_\_\_\_\_\_\_\_\_\_\_\_\_\_\_\_\_\_\_\_\_

\newpage

\section{Interpretación de los datos recolectados}

\subsection{Encuestas --- Orientación a Scrum}

\begin{longtable}{|L{5cm}|L{5cm}|L{5cm}|}
\hline
\rowcolor[gray]{0.9}
\textbf{Hallazgo} & \textbf{Producto de usuario} & \textbf{Prioridad} \\
\hline
\endhead
\multicolumn{3}{|c|}{\textit{(Registrar aquí los hallazgos de las encuestas una vez aplicadas)}} \\
\hline
 & & \\
\hline
 & & \\
\hline
 & & \\
\hline
 & & \\
\hline
 & & \\
\hline
\end{longtable}

\subsection{Observación --- Orientación a Scrum}

\begin{longtable}{|L{4cm}|L{5cm}|L{6cm}|}
\hline
\rowcolor[gray]{0.9}
\textbf{Observación} & \textbf{Criterio de aceptación} & \textbf{Definición de terminado} \\
\hline
\endhead
\multicolumn{3}{|c|}{\textit{(Registrar aquí los hallazgos de la observación una vez aplicada)}} \\
\hline
 & & \\
\hline
 & & \\
\hline
 & & \\
\hline
 & & \\
\hline
\end{longtable}

\subsection{Entrevista --- Orientación a Scrum}

\begin{longtable}{|L{4cm}|L{5cm}|L{6cm}|}
\hline
\rowcolor[gray]{0.9}
\textbf{Frase textual importante} & \textbf{Interpretación ágil} & \textbf{Producto propuesto} \\
\hline
\endhead
\multicolumn{3}{|c|}{\textit{(Registrar aquí las frases textuales de las entrevistas una vez aplicadas)}} \\
\hline
 & & \\
\hline
 & & \\
\hline
 & & \\
\hline
 & & \\
\hline
\end{longtable}

\newpage

\section{Situación del proceso actual}

\subsection{Identificación de actores (Swimlanes)}

A continuación, se identifican los actores que participan en los procesos actuales de retenciones y movimientos de caja en MINCH SRL.:

\begin{itemize}
    \item \textbf{Auxiliar Contable:} Encargado del registro diario de movimientos de caja, emisión de recibos de retención y mantenimiento de archivos Excel.

    \item \textbf{Contador General:} Responsable de la verificación de cálculos de retenciones, conciliación bancaria y generación de reportes financieros.

    \item \textbf{Gerente / Administrador:} Supervisor del área contable, toma decisiones basadas en los reportes financieros y autoriza movimientos.

    \item \textbf{Proveedor:} Persona o entidad que recibe el recibo de retención al momento de realizar una venta de bienes o servicios a la empresa.

    \item \textbf{Impuestos Nacionales (SIN):} Entidad gubernamental receptora de los pagos tributarios generados por las retenciones.
\end{itemize}

\subsection{Descripción del flujo actual --- Retenciones}

El proceso actual de retenciones se describe a continuación:

\begin{enumerate}
    \item \textbf{Disparador (Trigger):} El proveedor entrega una factura o documento de venta sin NIT al área contable.

    \item El Auxiliar Contable recibe la factura y verifica el tipo de bien o servicio adquirido (Servicios o Bienes).

    \item El Auxiliar abre el archivo Excel de retenciones y busca el último número correlativo registrado.

    \item Calcula manualmente el monto de la retención aplicando los porcentajes correspondientes según el tipo (RC-IVA 13\% + IT 3\% para Servicios; IUE 5\% + IT 3\% para Bienes).

    \item Registra los datos en el Excel: fecha, proveedor, monto, tipo, número de retención.

    \item Elabora manualmente el recibo de retención en un formato Word o Excel.

    \item El Contador General revisa y firma el recibo de retención.

    \item Se entrega el recibo al proveedor y se archiva una copia física.

    \item \textbf{Punto de dolor:} El cálculo manual es propenso a errores. La numeración correlativa puede duplicarse si dos personas registran al mismo tiempo. La búsqueda de registros históricos es lenta.
\end{enumerate}

\subsection{Descripción del flujo actual --- Movimientos de caja}

\begin{enumerate}
    \item \textbf{Disparador (Trigger):} Se realiza un ingreso o egreso en una de las cuentas bancarias de la empresa.

    \item El Auxiliar Contable recibe el comprobante de la transacción (transferencia, cheque, depósito).

    \item Abre el archivo Excel correspondiente a la cuenta bancaria (BNB, BCP, BISA, etc.).

    \item Identifica la última fila registrada y añade los datos del nuevo movimiento: fecha, descripción, monto, tipo (débito/crédito).

    \item Calcula manualmente el saldo acumulado sumando o restando el monto al saldo anterior.

    \item Si corresponde, emite un recibo de pago en formato Word.

    \item Periódicamente, el Contador General realiza la conciliación comparando el Excel con el extracto bancario.

    \item \textbf{Punto de dolor:} El registro manual consume aproximadamente el 40\% de la jornada laboral. Las versiones del Excel se duplican al compartirlas por correo. Los errores de cálculo de saldos son frecuentes. La generación de reportes mensuales requiere consolidar manualmente varias hojas.
\end{enumerate}

\subsection{Puntos de dolor identificados}

\begin{longtable}{|L{3cm}|L{4cm}|L{4cm}|L{4cm}|}
\hline
\rowcolor[gray]{0.9}
\textbf{Área} & \textbf{Punto de dolor} & \textbf{Impacto} & \textbf{Posible solución} \\
\hline
\endhead
Retenciones & Cálculo manual de descuentos propenso a errores & Alto & Automatizar cálculo de retenciones con tasas configurables \\
\hline
Retenciones & Duplicación de numeración correlativa & Alto & Numeración automática con bloqueo concurrente \\
\hline
Retenciones & Búsqueda lenta de registros históricos & Medio & Base de datos centralizada con búsqueda \\
\hline
Caja & Registro manual en Excel (40\% del tiempo) & Alto & Formulario web de registro rápido \\
\hline
Caja & Versiones contradictorias del Excel & Alto & Base de datos centralizada con control de acceso \\
\hline
Caja & Errores en cálculo de saldos acumulados & Alto & Cálculo automático de saldos con window functions \\
\hline
Caja & Conciliación bancaria lenta y manual & Medio & Reportes automáticos con saldos por cuenta \\
\hline
Reportes & Generación manual de reportes mensuales & Medio & Exportación automática a PDF y Excel \\
\hline
General & Dependencia de herramientas externas (Excel, Word) & Alto & Sistema web unificado \\
\hline
\end{longtable}

\end{document}
